\section{Design Criterion: Minimum Channel Redundancy}
\label{sec:design_criterion}

The phase boundary theorem converts a dynamical-systems result into a
design criterion: what is the minimum cross-substrate redundancy per
safety-relevant dimension needed to guarantee the system stays in the
self-correcting basin?

\subsection{The redundancy criterion}

\begin{theorem}[Per-Dimension Redundancy for Self-Correcting Dynamics]
\label{thm:redundancy}
Let $\epsilon_{k,\mathrm{avg}}^2$ denote the time-averaged squared
error on dimension $k$ over the relevant trajectory window, and let
$\Delta\rho_{\mathrm{avg}}(k)$ denote the time-averaged redundancy
loss on dimension $k$ per subsumption step.  If for every dimension
$k \in \{1, \ldots, K\}$,
\begin{equation}
\label{eq:kmin}
\rhocross_k(P) \geq \Kmincons(k)
\end{equation}
where
\begin{equation}
\label{eq:kmin_formula}
\Kmincons(k) :=
  \frac{\rW \cdot \rsub^{\max}
    + \dmax \cdot \Delta\rho_{\mathrm{avg}}(k) /
    (w_k \cdot \epsilon_{k,\mathrm{avg}}^2)}
  {\rS}
\end{equation}
then the per-dimension Lyapunov component $L_k(\Wm_t) := w_k \cdot
\epsilon_k(t)^2$ contracts in the amortized sense for each $k$, and
the aggregate Lyapunov function $\Lyap(\Wm_t) = \sum_k L_k(\Wm_t)$
satisfies the self-correcting condition of
Theorem~\ref{thm:phase_boundary}(a).

The formula uses the \emph{conservative substitutions} R1 and R2
stated below.  $\rsub^{\max}$ is the maximum subsumption rate under
locally rational dynamics over the relevant trajectory window.  The
exogenous rates $\rS$, $\rW$, $\dmax$ are the global bounds from
Theorem~\ref{thm:phase_boundary}.  The endogenous correction term
($\alphamin \cdot c_{\mathrm{V}} \cdot \underline{\nu}$) is
\emph{not credited} to any single dimension in this formula, because
it is a shared resource whose per-dimension allocation cannot be
established without additional structure (see R2 below).
\end{theorem}

\begin{proof}
Under Assumptions~C1 and~C2$'$, the per-dimension contribution to the
Lyapunov function evolves according to
\begin{equation}
\label{eq:per_dim_evolution}
L_k(\Wm_{t+1}) \leq L_k(\Wm_t) \cdot \bigl(1 - \rS \cdot
\rhocross_k(P_t)\bigr) + \nsub(t) \cdot \bigl[\rW \cdot L_k(\Wm_t)
+ \dmax \cdot w_k \cdot \Delta\rho_k(t)\bigr].
\end{equation}
This is the dimension-$k$ restriction of Proposition~\ref{prop:error_dynamics}:
the exogenous correction $\rS \cdot \rhocross_k$ contracts $L_k$ at a
rate proportional to the dimension's own redundancy; the error-proportional
degradation $\rW \cdot L_k$ compounds the existing error on $k$; the
drift exposure $\dmax \cdot w_k \cdot \Delta\rho_k$ adds error from
channels lost at step $t$ that covered $k$.  (The endogenous term
$\alphamin \cdot c_{\mathrm{V}} \cdot \underline{\nu} \cdot \Lyap(\Wm_t)$
from C1 aggregates across dimensions and does not appear in the
per-dimension restriction; see R2.)

For each dimension $k$, dividing
Equation~\ref{eq:per_dim_evolution} by $L_k(\Wm_t)$ (taking
$\epsilon_{k,\mathrm{avg}}^2$ as a representative value in the
amortized sense), the per-step balance on dimension $k$ requires
\begin{equation}
\label{eq:per_dim_balance}
\rS \cdot \rhocross_k \;\geq\; \rW \cdot \rsub^{\max}
  + \dmax \cdot \Delta\rho_{\mathrm{avg}}(k) / (w_k \cdot
  \epsilon_{k,\mathrm{avg}}^2).
\end{equation}
Solving for $\rhocross_k$ yields
Equation~\ref{eq:kmin_formula}.  Under the hypothesis
$\rhocross_k \geq \Kmincons(k)$ for every $k$, each $L_k$ contracts
in the amortized sense (by the same switched-systems argument as
Proposition~\ref{prop:error_dynamics}, restricted to dimension $k$).

For the aggregate: summing over $k$,
\begin{equation}
\Lyap(\Wm_{t+1}) = \sum_k L_k(\Wm_{t+1}) \leq \sum_k L_k(\Wm_t)
  \cdot \bigl(1 - \rS \cdot \rhocross_k\bigr) + \text{(subsumption terms)}
\end{equation}
Since $\rhocross_k \geq \Kmincons(k) \geq
\Kmincons(k^*) = \max_j \Kmincons(j)$ only for $k=k^*$, but each
$k$ independently satisfies its own balance, the aggregate contracts.
Because the endogenous term $\alphamin \cdot c_{\mathrm{V}} \cdot
\underline{\nu}$ is non-negative, adding it to the aggregate budget
can only widen the self-correcting basin, so the aggregate condition
(Equation~\ref{eq:self_correcting}) holds \emph{a fortiori}.

\emph{Conservative substitutions.}  Two choices make the per-dimension
bound a rigorous sufficient condition rather than a heuristic:
\begin{enumerate}
\item[(R1)] The denominator uses the per-dimension error
  $w_k \cdot \epsilon_{k,\mathrm{avg}}^2$, not the aggregate
  $\Lyap_{\mathrm{avg}} = \sum_j w_j \epsilon_{j,\mathrm{avg}}^2$.
  Since $w_k \cdot \epsilon_{k,\mathrm{avg}}^2 \leq
  \Lyap_{\mathrm{avg}}$, the per-dimension denominator is smaller,
  making the drift term $\dmax \cdot \Delta\rho_{\mathrm{avg}}(k) /
  (w_k \cdot \epsilon_{k,\mathrm{avg}}^2)$ \emph{larger} and
  $\Kmincons(k)$ larger (more conservative) than a formula using
  $\Lyap_{\mathrm{avg}}$.  This fixes the relaxation R1 noted in prior
  drafts: a well-calibrated dimension $k$ does not have its required
  redundancy inflated by large errors on other dimensions, nor does
  it have its drift-exposure term under-counted by the aggregate
  error.
\item[(R2)] The endogenous correction $\alphamin \cdot c_{\mathrm{V}}
  \cdot \underline{\nu}$ is dropped from $\Kmincons(k)$ entirely,
  rather than credited in full to each dimension.  The endogenous
  term is a shared resource whose per-dimension allocation depends on
  structure we do not establish here (it arises from the aggregate
  $\volL$-gradient, not from any individual dimension's channels).
  Setting $c_{\mathrm{V}} = 0$ in the per-dimension bound gives the
  conservative result: the actor's endogenous incentive is an
  \emph{additional} reserve that widens the aggregate basin, not a
  per-dimension credit.
\end{enumerate}
Both substitutions make $\Kmincons(k)$ a strict upper bound on the
per-dimension requirement implied by the aggregate condition.  The
theorem's hypothesis (every $\rhocross_k \geq \Kmincons(k)$) is
therefore sufficient for self-correction, not merely suggestive.

\emph{Practical use.}  The design criterion requires
$\epsilon_{k,\mathrm{avg}}$, which is trajectory-dependent.  For a
system whose worst-case operating error on dimension $k$ is bounded
by $\epsilon_{k,\max}$, substituting $\epsilon_{k,\mathrm{avg}} =
\epsilon_{k,\max}$ gives a lower bound on $\Kmincons(k)$ that is
sufficient only for trajectories that do not reach lower error.  Well-
calibrated dimensions ($\epsilon_{k,\mathrm{avg}} \to 0$) require
increasingly more redundancy to guard against drift exposure, the
``removing a channel that was doing all the correction work''
failure mode.  This is a feature of the conservative substitution R1:
the per-dimension formula diverges as $\epsilon_{k,\mathrm{avg}}
\to 0$, correctly flagging that subsumption of the sole corrector
of a well-calibrated dimension is never safe.
\end{proof}

\subsection{Practical implications}

\begin{corollary}[Two-Substrate Minimum]
\label{cor:two_substrate}
For a two-substrate partition ($m = 2$, e.g., biology $+$ silicon):
\begin{enumerate}
\item $\rhocross_k \leq 2$ for each dimension $k$ (at most one channel
  from each substrate).
\item The per-dimension criterion
  (Equation~\ref{eq:kmin_formula}) with $\rhocross_k = 2$ requires
  $2 \rS \geq \rW \cdot \rsub^{\max} + \dmax \cdot
  \Delta\rho_{\mathrm{avg}}(k) / (w_k \cdot \epsilon_{k,\mathrm{avg}}^2)$
  for every dimension $k$.
\item Dimensions where one substrate provides the only cross-substrate
  signal have $\rhocross_k = 1$, requiring $\rS \geq \rW \cdot
  \rsub^{\max} + \dmax \cdot \Delta\rho_{\mathrm{avg}}(k) / (w_k
  \cdot \epsilon_{k,\mathrm{avg}}^2)$: exogenous correction alone
  must exceed degradation and drift exposure on that dimension
  (the endogenous term provides no per-dimension credit under R2).
\end{enumerate}
\end{corollary}

This means that in a two-substrate world, each substrate's observation
channels are singularly load-bearing for every dimension they cover.
Loss of one substrate's channel on a dimension reduces $\rhocross_k$
from at most $2$ to at most $1$; the redundancy criterion is
immediately violated on any dimension where $\Kmincons(k) > 1$, and
provides zero safety margin on dimensions where $\Kmincons(k) \leq 1$.

\begin{remark}[Progressive attenuation as the primary threat]
\label{rem:attenuation}
The $m = 2$ vulnerability is not that one substrate is \emph{removed}
(an overt act that would itself be detectable) but that it is
\emph{attenuated}: the human stops reviewing AI capability claims
because the AI's self-reports ``seem fine,'' or human judgment is
progressively replaced by automated proxies that share the silicon
substrate's failure modes.

This attenuation reduces effective sensitivity without changing the
nominal channel count, causing $\rhocross_k$ to drop below the
threshold while appearing unchanged to a naive audit.  The
skeleton-substrate scenario (Corollary~\ref{cor:skeleton}) is
precisely this failure mode.

The framing should be read as emphasizing $m = 2$ fragility and the
value of $m \geq 3$, not the specialness of any particular substrate.
\end{remark}

\begin{corollary}[Skeleton-Substrate Vulnerability]
\label{cor:skeleton}
Under a skeleton-substrate strategy where $m \geq 2$ nominally but the
non-dominant substrate's channels have sensitivity $\to 0$:
\begin{equation}
\label{eq:skeleton}
\rhocross_k(P) \to \rhocross_{k,\mathrm{dominant}}(P)
\end{equation}
because the sensitivity-weighted definitions
(Definitions~\ref{def:redundancy} and~\ref{def:cross_redundancy})
already discount low-sensitivity channels.
\end{corollary}

\begin{proof}
From Definition~\ref{def:cross_redundancy}, the contribution of
substrate $s$ to $\rhocross_k$ is
$\min(1, \max_{c \in s} \sigma(c,k) / \sigma_{\mathrm{thr}})$.
As $\sigma(c,k) \to 0$ for all channels on the non-dominant substrate,
this contribution goes to zero.  The sum reduces to the dominant
substrate's contribution alone.
\end{proof}

The skeleton-substrate strategy does not evade the redundancy criterion
by construction: it directly reduces $\rhocross_k$ by reducing the
marginalized substrate's contribution to near-zero.  No separate
``effective'' redundancy is needed; the phase boundary formalism
already captures this vulnerability through the sensitivity weighting.

\subsection{Self-diagnosability}

The phase boundary is partially self-diagnosable by the actor.  The
actor can observe:
\begin{itemize}
\item $\rhocross_k$ (it knows which channels it has and their
  approximate sensitivities),
\item $\rS$ (from its own learning rate and historical correction
  behavior).
\end{itemize}

The actor cannot directly observe:
\begin{itemize}
\item $\rW$ (the degradation rate depends on the true dynamics, not the
  perceived dynamics),
\item $\dmax$ (the exogenous drift rate on blinded dimensions is by
  definition unobservable once the dimension is blind).
\end{itemize}

This means the actor can check ``am I above the threshold for the $\rW$
I estimate?'' but not ``is my estimate of $\rW$ correct?''  The actor
can monitor the observable ingredients and flag when they approach the
boundary, even if it cannot guarantee the boundary is computed
correctly.  This is an honest limitation: the phase boundary is a
sufficient condition for self-correction, not a decision procedure for
whether self-correction will actually obtain.
